Sei $F$ die Antriebskraft, $c_\mathrm{R}$ der Faktor des Rollwiderstands, $m_1$ die Masse des schleppenden und $m_2$ die des abgeschleppten Autos.

\begin{abcliste}
\abc
Beide Autos zusammen als ein System: Die Antriebskraft gegen beide Rollwiderstände beschleunigt die gesamte Masse. Die Seilkraft ist eine innere Kraft dieses Systems und kommt nicht vor.
\begin{align}
 a &= \aF \\
   &= \frac{\FA-\cR\cdot(\ma+\mb)\cdot\g}{\ma+\mb} \\
   &= \a \approx \result{\aP{0}{2}}
\end{align}

\abc
Das abgeschleppte Auto allein: Das Seil zieht es nach vorn, sein Rollwiderstand hält es zurück.
\begin{align}
 \ssc{F}{S} - c_\mathrm{R}\,m_2\,g &= m_2\,a \\
 \ssc{F}{S} &= \FSF \\
   &= \mb\cdot\a+\cR\cdot\mb\cdot\g \\
   &= \FS \approx \result{\FSP{0}{2}}
\end{align}
Das Seil gibt nur weiter, was das abgeschleppte Auto braucht.
\end{abcliste}
